\section{Logic, Sets and Numbers}

\subsection{Logic}

\begin{definition}
    A \textbf{mathematical Statement} is a well-formulated assertion that is strictly either true or false.
\end{definition}

\begin{example}
    "Every prime number is odd." is a false statement.
\end{example}

\begin{definition}
A \textbf{Statement} that depends on one or multiple variables is called a \textbf{predicate} and denoted by $P(x), Q(x), \dots$.
\end{definition}

\begin{example}
$$P(x) \coloneqq x > 2$$
\end{example}

\begin{definition}
    The negation of a mathematical Statement $A$ is defined as $\neg A \coloneqq A \text{ is false}$.
\end{definition}

\begin{example}
    $A \coloneqq \forall x \in \mathbb{R}: x > 2$
    $\neg A \coloneqq \exists x \in \mathbb{R}: x \leq 2$
\end{example}

\begin{definition}
    \textbf{Quantifiers} are short hand notation for making statements about predicates.
    $$\forall x \in \mathbb{R}: P(x) \quad \iff \quad \text{for all} x \in \mathbb{R} \text{ holds } P(x)$$
    $$\exists x \in \mathbb{R}: P(x) \quad \iff \quad \text{there exists} x \in \mathbb{R} \text{ s.t. } P(x)$$
    $$\exists ! x \in \mathbb{R}: P(x) \quad \iff \quad \text{there exists exactly one} x \in \mathbb{R} \text{ s.t. } P(x)$$
    $$\nexists x \in \mathbb{R}: P(x) \quad \iff \quad \text{there exists no} x \in \mathbb{R} \text{ s.t. } P(x)$$
\end{definition}

\begin{example}
    $$\forall x \in \mathbb{R}: x^2 \geq 0$$
    $$\forall x \in \mathbb{R} \exists y \in \mathbb{R}: x = y^2$$
\end{example}

\begin{important}
    The oderder of quantifiers matters. $\forall x \exists y$ is not equivalent to $\exists y \forall x$.
\end{important}

\begin{definition}
    Basic \textbf{logical connectives} are defined as:
    $$A \land B \coloneqq A \text{ and } B \text{ are true}$$
    $$A \lor B \coloneqq A \text{ or } B \text{ is true}$$
    $$A \implies B \coloneqq A \text{ implies } B \text{ is true}$$
    $$A \iff B \coloneqq A \text{ is equivalent to } B \text{ is true}$$
\end{definition}

\begin{remark}
    Observe that equivalence holds iff. $A \implies B \land B \implies A$
\end{remark}

\begin{lemma}
    From the above definition it follows that:
    $$ \neg A \lor B \iff A \implies B$$
    This can be shown by evaluating the truth table for both statements.
\end{lemma}

\begin{theorem}
    The following equivalences hold:
    $$\neg (\forall x P(x)) \iff \exists x \neg P(x)$$
    $$\neg (\exists x P(x)) \iff \forall x \neg P(x)$$
    aswell as
    $$\forall x (P(x) \land Q(x)) \iff (\forall x P(x)) \land (\forall x Q(x))$$
    $$\exists x (P(x) \lor Q(x)) \iff (\exists x P(x)) \lor (\exists x Q(x))$$
\end{theorem}

\subsection{Sets}

\begin{definition}
    A \textbf{set} is a collection of elements, that are listed or characterized by a certain property. 
    The empty set is denoted by $\emptyset$ and is the only set that contains no elements.
\end{definition}

\begin{definition}
    An element $x$ is \textbf{contained in a set} $A$ iff. $x \in A$ else it is not contained in $A$ denoted by $x \notin A$.
\end{definition}

\begin{definition}
    Let $A$ and $B$ be sets. $A$ is a \textbf{subset} of $B$ denoted by $A \subset B$ iff. 
    $$\forall x (x \in A \implies x \in B)$$
    $A$ is a \textbf{proper subset} of $B$ denoted by $A \subsetneq B$ iff. 
    $$A \subset B \land A \neq B$$
\end{definition}

\begin{definition}
    Let $A$ and $B$ be sets.
    \begin{itemize}
    \item The \textbf{union} of $A$ and $B$ is defined as $A \cup B \coloneqq \{x \mid x \in A \lor x \in B\}$.
    \item The \textbf{intersection} of $A$ and $B$ is defined as $A \cap B \coloneqq \{x \mid x \in A \land x \in B\}$.
    \item The \textbf{difference} of $A$ and $B$ is defined as $A \setminus B \coloneqq \{x \mid x \in A \land x \notin B\}$.
    \item The \textbf{Cartesian product} of $A$ and $B$ is defined as $A \times B \coloneqq \{(a, b) \mid a \in A \land b \in B\}$.
\end{itemize}
\end{definition}

\begin{remark}
    Note that if $B \subset A$ then the diffrence $A \setminus B$ is called the complement of $B$ in $A$ and is denoted by $B^c$.
\end{remark}

\subsection{Numbers}

The basic and often used sets of numbers are:

\begin{itemize}
    \item $\mathbb{N}$: The set of natural numbers 1, 2, 3, 4, ...
    \item $\mathbb{Z}$: The set of integers ... -2, -1, 0, 1, 2, ...
    \item $\mathbb{Q}$: The set of rational numbers $\frac{p}{q}$ where $p, q \in \mathbb{Z}$
    \item $\mathbb{R}$: The set of real numbers
\end{itemize}

\begin{remark}
    Some important subsets of the real numbers are $\mathbb{R}^{+}_{0} / \mathbb{R}^{+}$ and $\mathbb{R}^{-}_{0} / \mathbb{R}^{-}$ which define the intervalls $[0, \infty) / (0, \infty)$ and $(-\infty, 0] / (-\infty, 0)$.
\end{remark}

\begin{definition}
    Continuous subsets are called \textbf{intervalls.} An intervall is called
    \begin{itemize}
        \item \textbf{open} if it does not contain its endpoints $I = (a, b) \coloneqq \{x \mid a \lt x \lt b\}$
        \item \textbf{closed} if it contains its endpoints $I = [a, b] \coloneqq \{x \mid a \leq x \leq b\}$
        \item \textbf{half-open} if it contains one endpoint and not the other $I = [a, b) \coloneqq \{x \mid a \leq x \lt b\}$
    \end{itemize}
\end{definition}

\begin{remark}
    An arbitrary subset of $\mathbb{R}$ is called open if it is a union of open intervalls and closed if its complement is open.
\end{remark}

\begin{definition}
    An \textbf{upper (lower) bound} of a subset $X \subset \mathbb{R}$ is an element $y \in \mathbb{R}$ such that:
    $$\forall x \in X: x \leq (\geq) y$$
    The \textbf{maximum (minimum)} of a subset $X \subset \mathbb{R}$ is the element $x_{0} \in X$ such that:
    $$\forall x \in X: x \leq (\geq) x_{0}$$
    Often denoted as \textbf{max(X)} and \textbf{min(X)}.
\end{definition}

\begin{definition}
    An intervall $I = (a, b)$ is \textbf{bounded} if there exist upper and lower bounds.
\end{definition}

\begin{remark}
    A bounded and closed intervall is often refered to as \textbf{compact}.
\end{remark}
\begin{definition}
    The \textbf{Supremum} $sup(X)$ of a subset $X \subset \mathbb{R}$ is the least upper bound of $X$. and can be characterized by:
    $$\forall x \in X: x \leq sup(X) \land \\ \forall \epsilon \gt 0: \exists x \in X: x \gt sup(X) - \epsilon$$
\end{definition}

\begin{definition}
    The \textbf{Infimum} $inf(X)$ of a subset $X \subset \mathbb{R}$ is the greatest lower bound of $X$. and can be characterized by:
    $$\forall x \in X: x \geq inf(X) \land \\ \forall \epsilon \gt 0: \exists x \in X: x \lt inf(X) + \epsilon$$
\end{definition}

\begin{example}
    Let $X = [-2, 1)$ we have a minimum of $x_0 = -1$ but no maximum. Observe that it is bounded from above by $1$ and $2$ or any number larger than $1$ but only $1$ is the least upper bound (Supremum).
\end{example}

\begin{theorem}
    The \textbf{maximum and minimum are unique} aslong as they exist.
\end{theorem}

\begin{proof}
    Assume there exists $x_1, x_2 \in X$ such that $x_1 = max(X)$ and $x_2 = max(X)$. As both $x_1$ and $x_2$ are maximum of $X$ we have
    $$x_1 \geq x_2 \land x_2 \geq x_1 \implies x_1 = x_2$$
\end{proof}

\begin{definition}
    \textbf{Axioms of Addition}
    \begin{itemize}
        \item Commutativity: $x + y = y + x$
        \item Associativity: $x + (y + z) = (x + y) + z$
        \item Identity element: There exists an element $e \in X$ such that $x + e = x$ for all $x \in X$.
        \item Inverse element: For every $x \in X$ there exists an element $x^{-1} \in X$ such that $x + x^{-1} = e$.
    \end{itemize}
\end{definition}

\begin{remark}
    A set $X$ together with an operation $+$ that satisfies the axioms of addition is called a \textbf{abelian group}.
\end{remark}

\begin{example}
    $(\mathbb{Z}, +)$ is an abelian group.
\end{example}

\begin{definition}
    \textbf{Axioms of Multiplication}
    \begin{itemize}
        \item Commutativity: $x \cdot y = y \cdot x$
        \item Associativity: $x \cdot (y \cdot z) = (x \cdot y) \cdot z$
        \item Identity element: There exists an element $e \in X$ such that $x \cdot e = x$ for all $x \in X$.
        \item Inverse element: For every $x \in X$ there exists an element $x^{-1} \in X$ such that $x \cdot x^{-1} = e$.
    \end{itemize}
\end{definition}

\begin{definition}
    \textbf{Distributivity} of multiplication over addition:
    $$\forall x, y, z \in X: x \cdot (y + z) = x \cdot y + x \cdot z$$
\end{definition}

\begin{remark}
    A set $X$ together with operations $+$ and $\cdot$ that satisfies the axioms of addition and multiplication and distributivity is called a \textbf{field}.
\end{remark}

\begin{definition}
    \textbf{Order Axioms}
    \begin{itemize}
        \item reflexivity: $\forall x \in X: x \leq x$
        \item antisymmetry: $\forall x, y \in X: x \leq y \land y \leq x \implies x = y$
        \item transitivity: $\forall x, y, z \in X: x \leq y \land y \leq z \implies x \leq z$
        \item totality: $\forall x, y \in X: x \leq y \lor y \leq x$
    \end{itemize}
\end{definition}

\begin{definition}
    A set $X$ is said to be \textbf{(order)complete} iff for every non empty subsets $A, B \subseteq X$
    $$\forall a \in A \forall b \in B: a \leq b$$
    $$\implies \exists c \in X \quad \text{s.t.} \quad \forall a \in A: a \leq c \land \forall b \in B: c \leq b$$
\end{definition}

\begin{example}
    The real numbers $\mathbb{R}$ satisfy axiomatic properties for addition, multiplication, distributivity and order. Furthermore $\mathbb{R}$ is an order complete field.
\end{example}

\begin{example}
    $\mathbb{Q}$ is not order complete. Let $A = \{x \in \mathbb{Q} \mid x^2 \leq 2\}$ and $B = \{x \in \mathbb{Q} \mid x^2 \geq 2\}$. Then $\forall a \in A \forall b \in B: a \leq b$ but there is no $c \in \mathbb{Q}$ such that $c^2 = 2$.
\end{example}

\begin{theorem}
    Every \textbf{non empty, bounded from above (below)} subset $X \subset \mathbb{R}$ has a \textbf{supremum (infimum)}.
\end{theorem}

\begin{proof}
    Let $B$ be the set of all upper bounds of $X$. As $X$ is bounded from above and non empty, $B$ is non empty. Therfore
    $$\forall x \in X: x \leq b \quad \forall b \in B$$
    By the order completenes of $\mathbb{R}$ there exists a $c \in \mathbb{R}$ such that
    $$\forall x \in X: x \leq c \land \forall b \in B: c \leq b$$
    Hence $c$ is the least upper bound of $X$ and therefore $c = sup(X)$. Assume there exists $c' \in \mathbb{R}$ such that $c' = sup(X)$ and $c' \neq c$. Then by definition of $sup(X)$ we have
    $$c \leq c' \land c' \leq c \implies c = c'$$
\end{proof}

\begin{theorem}
The \textbf{Archimedian Principle} says that
\begin{itemize}
    \item For $x \in \mathbb{R}$ and $y \gt 0$ there exists an element $n \in \mathbb{N}$ such that $n * y \gt x$.
    \item For every $\epsilon \gt 0$ there exists an element $n \in \mathbb{N}$ such that $\frac{1}{n} \lt \epsilon$.
\end{itemize}
\end{theorem}

\begin{proof}
    We prove the first theorem by contradition. Assume it does not hold hence there exists $x \in \mathbb{R}$ and $y \gt 0$ such that
    $$\forall n \in \mathbb{N}: n * y \leq x$$
    This implies that $x$ is an upper bound for the set $X = \{n * y \mid n \in \mathbb{N}\}$. As $X$ is non empty and bounded from above, it has a supremum $c = sup(X)$. As $y \gt 0$ $c - y$ can not be an upper bound of $X$ hence there exists $n_0 \in \mathbb{N}$ such that
    $$c - y \lt n_0 * y \implies c \lt (n_0 + 1) * y$$
    This contradicts the assumption that $c$ is the supremum of $X$.
\end{proof}    
\begin{proof}
    The second theorem is a direct consequence of the first theorem. It can be shown by setting $x = 1$ in the first theorem.
\end{proof}    
    

\begin{theorem}
The Young Inequality states that for all $\epsilon \gt 0$ and all $x, y \in \mathbb{R}$:
$$2\lvert xy \rvert \leq \epsilon x^2 + \frac{1}{\epsilon} y^2$$
\end{theorem}

\begin{proof}
    $$\left( \sqrt{\epsilon}x - \frac{1}{\sqrt{\epsilon}}y \right)^2 \geq 0 \implies \epsilon x^2 - 2xy + \frac{1}{\epsilon} y^2 \geq 0 \implies \epsilon x^2 + \frac{1}{\epsilon} y^2 \geq 2xy$$
\end{proof}

\subsubsection{Vectors and Vector Spaces}

\begin{definition}
    The \textbf{euclidian space} $\mathbb{R}^n$ is the set of all vectors $x = (x_1, x_2, \dots, x_n)$ where $x_i \in \mathbb{R}$.
\end{definition}

\begin{definition}
    \textbf{Addition} and \textbf{scalar multiplication} are defined as:
    $$x + y = (x_1 + y_1, x_2 + y_2, \dots, x_n + y_n)$$
    $$\lambda x = (\lambda x_1, \lambda x_2, \dots, \lambda x_n)$$
\end{definition}

\begin{definition}
    A \textbf{vector space} has to satisfy the following axioms:
    \begin{itemize}
        \item Commutativity: $x + y = y + x$
        \item Associativity: $x + (y + z) = (x + y) + z$
        \item Identity element: There exists an element $e \in X$ such that $x + e = x$ for all $x \in X$.
        \item Inverse element: For every $x \in X$ there exists an element $x^{-1} \in X$ such that $x + x^{-1} = e$.
        \item Distributivity: $x \cdot (y + z) = x \cdot y + x \cdot z$
    \end{itemize}
\end{definition}

\begin{definition}
    For $x, y \in \mathbb{R}^n$, the \textbf{standard scalar product} is defined as:
    $$x \cdot y = \langle x, y \rangle \coloneqq \sum_{i=1}^n x_i y_i$$
    The \textbf{euclidean norm} is defined as:
    $$||x|| \coloneqq \sqrt{\sum_{i=1}^n x_i^2}$$
    The \textbf{euclidean distance} between $x$ and $y$ is:
    $$d(x, y) \coloneqq \sqrt{\sum_{i=1}^n (x_i - y_i)^2}$$
\end{definition}

\begin{theorem}
    The triangle- and the cauch-schwarz inequality hold for all $x, y, z \in \mathbb{R}^n$:
    $$||x - z|| \leq ||x - y|| + ||y - z||$$
    $$|\langle x, y \rangle| \leq ||x|| ||y||$$
\end{theorem}

\subsubsection{Complex Numbers}

\begin{definition}
    The \textbf{imaginary unit} $i$ is defined as $i^2 = -1$.
\end{definition}

\begin{definition}
    The \textbf{complex numbers} $\mathbb{C}$ are the set of numbers
    $$\mathbb{C} = \{a + ib \mid a, b \in \mathbb{R}\}$$
    where $a$ is the \textbf{real part} $\text{Re}(z)$, and $b$ is the \textbf{imaginary part} $\text{Im}(z)$.
\end{definition}

\begin{definition}
    Arithmetic of complex numbers $z_1 = x_1 + i y_1$ and $z_2 = x_2 + i y_2$:
    \begin{itemize}
        \item Addition: $z_1 + z_2 = (x_1 + x_2) + i(y_1 + y_2)$
        \item Multiplication: $z_1 \cdot z_2 = (x_1 x_2 - y_1 y_2) + i(x_1 y_2 + y_1 x_2)$
    \end{itemize}
\end{definition}

\begin{definition}
    The \textbf{complex conjugate} of $z = x + iy$ is $\bar{z} = x - iy$.
\end{definition}

\begin{important}
    Division is calculated by expanding with the conjugate of the denominator:
    $$\frac{z}{w} = \frac{z \cdot \bar{w}}{w \cdot \bar{w}} = \frac{z \cdot \bar{w}}{|w|^2}$$
\end{important}

\begin{remark}
    Some usfull properties of the conjugate:
    \begin{itemize}
        \item $z \cdot \bar{z} = |z|^2$
        \item $z + \bar{z} = 2\text{Re}(z)$ and $z - \bar{z} = 2i\text{Im}(z)$
        \item $\overline{z+w} = \bar{z} + \bar{w}$ and $\overline{z \cdot w} = \bar{z} \cdot \bar{w}$
    \end{itemize}
\end{remark}

\begin{definition}
    Geometrically, complex numbers can be represented in the Gaussian plane. The \textbf{absolute value} of $z = a + ib$ is its \textbf{distance from the origin}:
    $$|z| = \sqrt{a^2 + b^2}$$
    The distance between $z_1$ and $z_2$ is $d = |z_1 - z_2|$. 
\end{definition}

\begin{remark}
    Observe that the triangle inequality holds in $\mathbb{C}$: $|z + w| \leq |z| + |w|$.
\end{remark}

\begin{theorem}
    The Euler's formula states that for $t \in \mathbb{R}$:
    $$e^{it} = \cos(t) + i\sin(t)$$
    This allows writing trigonometric functions using complex exponentials:
    $$\cos(t) = \frac{e^{it} + e^{-it}}{2}, \quad \sin(t) = \frac{e^{it} - e^{-it}}{2i}$$
\end{theorem}

\begin{remark}
    The \textbf{Intuition} behind eulers formula is we can use Taylor series to show that:
    $$e^x = 1 + x + \frac{x^2}{2!} + \frac{x^3}{3!} + \dots = \sum_{n=0}^\infty \frac{x^n}{n!}$$
    Using $x=it$ we get:
    $$e^{it} = \sum_{n=0}^\infty \frac{(it)^n}{n!} = \dots = \cos(t) + i\sin(t)$$
\end{remark}

Eulers formula leads us to a new representation of complex numbers: \textbf{the polarcoordinates}.

\begin{definition}
    For $z \in \mathbb{C}$, the polar-coordinates are defined as:
    $$z = r(\cos(t) + i\sin(t)) = r e^{it}$$
    where $r = |z|$ is the absolute value and $t = \arg(z)$ is the argument.
\end{definition}

\begin{remark}
    To convert from cartesian to polar coordinates:
    $$r = \sqrt{a^2 + b^2}, \quad t = \arctan\left(\frac{b}{a}\right)$$
    Note that the arctan function only returns values in $(-\frac{\pi}{2}, \frac{\pi}{2})$, so we have to adjust the angle based on the quadrant of the complex number.
\end{remark}

\begin{example}
    Multiplication and division in polar-coordinates is much easier than in cartesian coordinates:
    $$z_1 \cdot z_2 = r_1 e^{it_1} \cdot r_2 e^{it_2} = r_1 r_2 e^{i(t_1 + t_2)}$$
    $$\frac{z_1}{z_2} = \frac{r_1 e^{it_1}}{r_2 e^{it_2}} = \frac{r_1}{r_2} e^{i(t_1 - t_2)} $$
\end{example}

\begin{example}
    So is exponentiation or taking roots:
    $$z^n = (r e^{it})^n = r^n e^{int}$$
    $$\sqrt[n]{z} = \sqrt[n]{r} e^{i\frac{t+2\pi k}{n}}, \quad k \in \{0, 1, \dots, n-1\}$$
\end{example}

\begin{definition}
    The $n$-th roots of unity are the solutions to $z^n = 1$, which are given by:
    $$z_k = e^{i\frac{2\pi k}{n}}, \quad k \in \{0, 1, \dots, n-1\}$$
\end{definition}

\begin{theorem}
    The sum of the roots of unit is always zero.
    $$\sum_{k=0}^{n-1} e^{i\frac{2\pi k}{n}} = 0$$
\end{theorem}

\begin{theorem}
    The \textbf{fundamental theorem of algebra} states that every polynomial of degree $n \geq 1$ has exactly $n$ roots in $\mathbb{C}$ (counting multiplicity). These roots appear in konplex-conjugate pairs.
\end{theorem}
